\chapter{Security}

Security is a default path, not a checklist at the end.
Small gaps in input handling or permissions become large incidents.

\section{Input handling}
Validate and sanitize input at the boundary.
Never pass raw input to SQL or shell commands.

\section{Secrets}
Store secrets in a dedicated manager.
Rotate regularly and audit access.

\section{Dependencies}
Track third-party dependencies and patch regularly.
Use allowlists for critical integrations.

\section{Rate limiting}
Rate limits protect your system and reduce abuse.
Apply limits per user and per IP when needed.

\begin{warningbox}{Warning: avoid security by obscurity}
Hidden endpoints are not secure endpoints.
Enforce authentication and authorization everywhere.
\end{warningbox}

\section{Threat modeling}
Threat modeling identifies the highest risk areas.
Start with data flows and trust boundaries.
Prioritize fixes based on impact and likelihood.

\section{Least privilege}
Grant the minimal permissions necessary.
Use scoped roles for services and users.
Review permissions regularly to prevent privilege creep.

\section{Encryption in transit}
Use TLS for all service-to-service and client traffic.
Do not allow plaintext connections in production.
Rotate certificates regularly.

\section{Encryption at rest}
Encrypt sensitive data at rest.
Manage keys in a dedicated key management system.
Access to decryption keys should be tightly controlled.

\section{Security headers}
Set standard security headers for web traffic.
Use \code{Strict-Transport-Security}.\\
Use \code{Content-Security-Policy}.\\
Headers reduce common browser-based attacks.

\section{Dependency scanning}
Automate dependency scans in CI.
Prioritize patching critical vulnerabilities.
Track exceptions and enforce expiration dates.

\section{Secrets management}
Secrets should never be stored in code.
Use a secrets manager with audit logs.
Rotate secrets on a schedule and on employee departure.

\section{Audit logging}
Audit logs capture security-relevant events.
Store audit logs in immutable storage.
Access to audit logs should be limited and monitored.

\section{Data classification}
Classify data as public, internal, confidential, or restricted.
Controls should vary based on classification.
Document how each class should be stored and accessed.

\section{Common attack vectors}
Plan for common vectors:
\begin{itemize}[nosep]
  \item SQL injection
  \item cross-site scripting
  \item SSRF
  \item credential stuffing
\end{itemize}
Controls should map to these known risks.

\section{Example security controls}
\begin{tabularx}{\textwidth}{@{}lX@{}}
\toprule
Risk & Control \\
\midrule
SQL injection & parameterized queries \\
Credential stuffing & rate limiting and MFA \\
SSRF & egress allowlists \\
\bottomrule
\end{tabularx}

\section{Security testing}
Add security tests for critical flows.
Use static analysis and dependency scanning.
Penetration tests are useful for mature systems.

\section{Incident response}
Define how to respond to security incidents.
Include communication plans and regulatory notifications.
Practice incident response regularly.

\section{Key rotation}
Rotate keys on a schedule and when incidents occur.
Automate rotation to reduce human error.
Monitor for failed rotation attempts.

\section{Secure defaults}
Default settings should be secure.
Opt-in to risky behavior rather than opt-out.
Secure defaults reduce accidental exposure.

\section{Authentication hardening}
Harden authentication with MFA and strong password policies.
Detect credential stuffing and apply rate limits.
Log and alert on unusual login patterns.

\section{Authorization boundaries}
Authorization checks must enforce tenant boundaries.
Use centralized policy evaluation when rules become complex.
Authorization errors should not leak resource existence.

\section{Input validation}
Validation prevents injection and unexpected behavior.
Use allowlists for critical inputs.
Reject input that cannot be parsed safely.

\section{Secure logging}
Logs should exclude secrets and personal data.
Use structured logging and redaction.
Audit access to sensitive logs.

\section{Data retention and deletion}
Define data retention policies by data class.
Implement deletion workflows that truly remove data.
Retention policies support compliance and privacy.

\section{Security reviews}
High-risk changes should go through security review.
Review threat models and mitigations.
Security reviews reduce blind spots in critical flows.

\section{Example security checklist}
\begin{itemize}[nosep]
  \item MFA required for privileged roles
  \item secrets stored in a manager
  \item dependency scans in CI
  \item audit logs enabled
\end{itemize}

\section{Secure SDLC}
Security should be part of the development lifecycle.
Add security checks to CI and code review.
Security gates prevent risky changes from reaching production.

\section{Supply chain security}
Software supply chains are common attack vectors.
Verify dependencies and pin versions.
Use signed artifacts when possible.

\section{Container and host security}
Harden containers with minimal base images.
Run services as non-root users.
Apply host-level hardening and monitoring.

\section{Network segmentation}
Segment networks by trust level.
Limit east-west traffic to required services.
Segmentation reduces lateral movement during incidents.

\section{WAF and edge protections}
Web application firewalls can block common attacks.
Edge protections should include rate limiting and bot detection.
Use them as defense in depth, not as the only control.

\section{Secrets rotation plan}
Define who owns secret rotation and when it happens.
Rotation should be automated and logged.
Test rotation in staging before production.

\section{Example secret rotation}
\begin{lstlisting}[language=YAML]
rotation:
  secret: "db_password"
  cadence_days: 90
  owner: "platform"
\end{lstlisting}

\section{Data access auditing}
Audit access to sensitive data.
Access reviews should be periodic.
Auditability is key for compliance and incident response.

\section{CSRF protection}
CSRF protection is required for browser-based sessions.
Use anti-CSRF tokens and same-site cookies.
CSRF controls are part of the auth surface.

\section{SSRF protection}
SSRF attacks can reach internal services.
Use egress allowlists and validate URLs.
Block access to metadata endpoints.

\section{Input encoding}
Encode output for the correct context.
HTML, JSON, and SQL each require different handling.
Output encoding is as important as input validation.

\section{Vulnerability management}
Track vulnerabilities and patch timelines.
Define SLAs for critical vulnerabilities.
Report patch status to stakeholders.

\section{Incident response runbook}
\begin{lstlisting}[language=YAML]
incident:
  detect: "alert or report"
  triage: "confirm scope"
  contain: "revoke access"
  eradicate: "patch and rotate keys"
  recover: "restore services"
\end{lstlisting}

\section{Security monitoring}
Monitor for anomalous access and data exfiltration.
Use alerts for spikes in failed logins.
Security monitoring should be continuous.

\section{Data classification table}
\begin{tabularx}{\textwidth}{@{}lX@{}}
\toprule
Class & Example data \\
\midrule
Public & marketing pages \\
Internal & internal docs \\
Confidential & customer emails \\
Restricted & payment details \\
\bottomrule
\end{tabularx}

\section{IAM policy example}
\begin{lstlisting}[language=JSON]
{
  "action": "read",
  "resource": "billing_invoices",
  "conditions": {"tenant_id": "tenant_7"}
}
\end{lstlisting}

\section{Data minimization}
Collect only the data you need.
Minimization reduces risk and compliance burden.
Review data collection regularly.

\section{Secure deletion}
Deletion should remove data and backups when required.
Document how deletion propagates across systems.
Secure deletion is part of privacy compliance.

\section{Security training}
Engineers should receive periodic security training.
Training reduces common mistakes in code and configuration.
Security knowledge decays without reinforcement.

\section{Secure configuration}
Misconfiguration is a leading cause of incidents.
Use configuration validation and policy checks.
Keep defaults locked down in production.

\section{Access reviews}
Access reviews should be periodic and documented.
Remove access that is no longer needed.
Access reviews reduce long-term exposure.

\section{Security metrics}
Track a small set of security metrics:
\begin{itemize}[nosep]
  \item failed login rate
  \item privileged access changes
  \item secrets rotation age
\end{itemize}
Metrics make security posture measurable.

\section{Example security event}
\begin{lstlisting}[language=JSON]
{
  "event": "privileged_access_granted",
  "actor": "user_7",
  "scope": "admin",
  "timestamp": "2026-02-18T12:00:00Z"
}
\end{lstlisting}

\section{Data residency}
Some data must stay within specific regions.
Document residency requirements and enforce them.
Residency violations can be compliance incidents.

\section{Secure backups}
Backups should be encrypted and access controlled.
Test backup restores regularly.
Backups are part of your security surface.

\section{Key management}
Keys should be separated by environment and purpose.
Use a dedicated key management system.
Audit key access and rotation events.

\section{Security reviews for changes}
High-risk changes should require an explicit security review.
Track security review approvals in code review tools.
Security reviews reduce surprise regressions.

\section{Configuration drift}
Configuration drift can open security gaps.
Use configuration as code and policy checks.
Drift detection should alert on unauthorized changes.

\section{Secure defaults checklist}
\begin{itemize}[nosep]
  \item deny by default for network access
  \item MFA required for privileged roles
  \item encryption enabled at rest
\end{itemize}

\section{Secrets sprawl}
Secrets sprawl happens when credentials are copied across services.
Track secrets inventory and owners.
Reduce sprawl by centralizing secret access.

\section{Security posture reviews}
Perform periodic security posture reviews.
Review logging, access, and patch status.
Posture reviews surface slow-moving risks.

\section{Patch cadence}
Define a patch cadence for critical dependencies.
Automate updates where possible.
Delayed patches are a leading incident cause.

\section{Session management}
Session tokens should be short-lived and rotated on privilege changes.
Invalidate sessions on logout and security events.
Session design is part of the authentication surface.

\section{Password storage}
Store passwords with a strong adaptive hash and unique salts.
Never log or transmit passwords in plaintext.
Password storage errors are permanent incidents.

\section{Token handling}
Treat access tokens as secrets with strict scope.
Use short expiry and refresh tokens for long-lived sessions.
Log token usage anomalies and revoke quickly.

\section{File upload security}
Validate file types and sizes at the edge.
Store uploads in isolated buckets with least privilege access.
Scan uploads for malware before processing.

\section{CORS policy}
CORS should be explicit and minimal.
Avoid wildcard origins for authenticated endpoints.
Document allowed origins and review regularly.

\section{Third-party risk}
Vendors can introduce security and compliance risk.
Review vendor security posture and data handling.
Terminate access when contracts end.

\section{Secrets in CI}
CI systems are a high-risk secrets target.
Use short-lived tokens and restrict build permissions.
Rotate build credentials regularly.

\section{Security champions}
Security champions create local expertise in teams.
They review changes and raise security issues early.
Champions improve the speed of security fixes.

\section{Security SLAs}
Define response SLAs for security findings.
Critical issues should have fast remediation targets.
SLAs ensure security work does not stall.

\section{Audit readiness}
Prepare evidence for audits continuously.
Keep policy docs, logs, and approvals organized.
Audit readiness reduces disruption during compliance reviews.

\begin{checklistbox}{Checklist: security}
\begin{itemize}[nosep]
  \item Validate all input.
  \item Rotate secrets and audit access.
  \item Patch dependencies and libraries.
\end{itemize}
\end{checklistbox}
